% TOP_PROBLEM: {"id": 7000022, "title": "Geometry of Curves and Surfaces — Problem 5.3", "category_id": 6, "category": {"id": 6, "name": "geometry", "display_name": "Geometry", "description": "Euclidean and non-Euclidean geometry, geometric structures.", "slug": "geometry", "order_index": 6, "created_at": "2026-07-31T15:26:25.671Z"}, "status": "open", "problem_number": "AMR-069-0022", "set_id": 13}
% FIRST_POSTED: 2026-10-01
% ENTRY_KIND: statement_only
% UnsolvedMath ID 7000022; source catalogue code AMR-069-0022
% !TeX program = lualatex
% Definition and sources only; this file is not an LLM solution attempt.
\documentclass[11pt,a4paper]{article}
\usepackage[margin=25mm]{geometry}
\usepackage{fontspec}
\setmainfont{lmroman10-regular.otf}[BoldFont=lmroman10-bold.otf,
 ItalicFont=lmroman10-italic.otf,BoldItalicFont=lmroman10-bolditalic.otf]
\usepackage{amsmath,amssymb,mathtools}
\usepackage{unicode-math}
\setmathfont{latinmodern-math.otf}
\AtBeginDocument{\let\setminus\smallsetminus}
\usepackage{xurl,hyperref}
\hypersetup{colorlinks=true,urlcolor=blue}
\setcounter{secnumdepth}{0}
\setlength{\parindent}{0pt}
\setlength{\parskip}{5pt}
\setlength{\emergencystretch}{3em}
\begin{document}
\section*{Geometry of Curves and Surfaces — Problem 5.3}\label{7000022}
{\small UnsolvedMath ID 7000022 · AMR-069-0022 · Geometry · AMR Open Problem Lists}
\subsection{Definitions and mathematical statement}
Let $\gamma:S^1\to\mathbb R^3$ be a closed rectifiable curve of
length $L>0$, and let $K=\operatorname{conv}(\gamma(S^1))$.
For a full-dimensional hull, let $A(K)$ be its boundary surface
area. Extend this quantity continuously to flat hulls: a planar
convex region has area counted twice, while a segment has area
zero. Equivalently, with $B$ the unit ball,
\[
 A(K)=\lim_{\varepsilon\downarrow0}
       \frac{V(K+\varepsilon B)-V(K)}{\varepsilon}.
\]
Here $+$ denotes Minkowski addition and $V$ is three-dimensional
volume.

The conjecture is
\[
                              A(K)\le\frac{L^2}{2\pi}.
\]
A planar circle of radius $L/(2\pi)$ attains the proposed value:
its hull is a disk counted on both sides. Determine also the
equality cases.

The curve is not required to be planar, simple, or contained in
the boundary of its convex hull; repeated traversal contributes
to $L$. In particular, a proof valid only when the curve splits
the hull boundary into two spanning disks would cover a
restricted case. The doubled-area convention for flat hulls is
essential to the stated constant and its extremal example.
\subsection{Short English statement}
Does a planar circle maximize the surface area of the convex hull
among all closed space curves of the same length, counting a
flat hull on both sides?
\subsection{Sources}
\begin{itemize}
\item[S1] ULAM.AI and the UnsolvedMath Contributors, supplied problems.json, record 7000022 (AMR-069-0022), AMR Open Problem Lists. Catalogue identity and source transcription; CC BY 4.0.
\url{https://huggingface.co/datasets/ulamai/UnsolvedMath}.
\item[S2] Ghomi - Open Problems in Geometry of Curves and Surfaces (2019), Problem 5.3, PDF page 14. Original source indexed by AMR-069-0022.
\url{https://people.math.gatech.edu/~ghomi/Papers/op.pdf}.
\end{itemize}
\end{document}
